\section{Ética e LGPD}

\begin{itemize}
  \item \textbf{Termo de consentimento (página inicial do formulário):}
  \begin{consentbox}
  \textit{\enquote{Esta pesquisa, de finalidade exclusivamente acadêmica, é conduzida por estudantes da disciplina de Medição e Experimentação em Engenharia de Software e investiga a percepção de desenvolvedores sobre velocidade e qualidade com assistentes de IA. Leva cerca de 4 minutos. Não coletamos nome, e-mail, empresa, IP nem localização, e as respostas são anônimas. A participação é voluntária e você pode abandonar o formulário a qualquer momento antes do envio; depois de enviada, a resposta não pode ser localizada nem excluída por ser anônima. Não há riscos previsíveis além de eventual desconforto leve. Contato: [e-mail do grupo]. Ao marcar a caixa abaixo, declaro ter 18 anos ou mais e concordar em participar.}}
  \end{consentbox}
  O avanço exige a caixa de aceite.
  \item \textbf{Anonimização:} sem login, sem coleta de identificadores ou metadados de conta; campos abertos orientam a não citar nomes. Os dados são tratados como anonimizados (LGPD, Art.\ 5\textsuperscript{o}, III e XI), caso em que a LGPD em regra não se aplica (Art.\ 12); o estudo adota voluntariamente finalidade, necessidade e minimização. Por ser pesquisa de opinião com participantes não identificados, aplica-se a dispensa de CEP/CONEP (Res.\ CNS 510/2016, Art.\ 1\textsuperscript{o}, parágrafo único, I).
  \item \textbf{Armazenamento:} respostas coletadas no Google Forms (servidores de terceiros, possivelmente no exterior; sem identificadores, portanto sem dado pessoal transferido). O CSV exportado fica em repositório privado no GitHub, com acesso restrito aos 5 integrantes. Ao fim da disciplina, os dados brutos são destruídos e só a versão anonimizada (textos livres revisados) é disponibilizada.
\end{itemize}
